% !TeX root = main.tex

\section{Related Work}
\label{sec:Related-Work}

This section collects reported works on the intersection of HPC and software engineering which is relevant to the role and the increased complexity of RSE activities.
As described recently by Grannan et al.~\cite{grannan2020understanding} HPC software projects have evolved to greatly take advantage of software engineering practices.
McInnes et al.~\cite{mcinnes2021community} propose embracing community software ecosystems so HPC projects can achieve their science goals through enriching the synergies with other communities with similar software needs.
In this regard, Bartlett et al.~\cite{10.14529/jsfi170104}, and Heroux et al.~\cite{heroux2020e4s} present the xSDK and E4S ecosystems, respectively, that provide community guidelines and policies towards a coordinated software interoperability, performance portability, and sustainability strategy.

Carver et al.~\cite{4222616} provide lessons learned from surveying five HPC applications and the gaps and opportunities for applying software engineering practices.
One important lesson is the difficulty of verification and validation.
Schmidberger and Brügge~\cite{6341492} conducted a survey among HPC projects that highlighted the trade-offs between software engineering practices and the time and resources required as one reason preventing further adoption.
Dubey et al.~\cite{dubey2014evolution} describes the symbiotic relationship between scientific research and good software engineering on the widely used FLASH code.
Lessons learned on the development of GPU applications are presented by van Werkhoven et al.~\cite{van2020lessons}.
In their work, they argue that GPU development comes with specific challenges and that code must be of sufficient quality.
In a recent paper, Pachev et al.~\cite{10.1145/3491418.3535124} present a CI framework for HPC applications.
Their work is based on GitHub Actions which cites similar advantages for the use of self-hosting CI runners as in this work.

At the beginning of the last decade, Schmidberger and Schmidberger~\cite{6341491} proposed the concept of ``software engineering as a service'' for HPC to lower existing barriers. Heroux et al.
~\cite{10.1007/978-3-030-44728-1_6} propose an iterative workflow, Productivity and Sustainability Improvement Planning (PSIP), that allows HPC software teams to identify development bottlenecks. Eisty et al.~\cite{10.1007/978-3-030-50436-6_33} used this methodology to develop a testing framework to test non-deterministic parallel research software.
Antonioletti et al.~\cite{antonioletti2000software}, anticipated the rapid evolution of software engineering practices for HPC codes as we entered the new century. More recently, Heroux~\cite{10078171} proposes the concept of research software science (RSS) to elevate the RSE practice and its impact due to the complexity of research software. 
Overall, the relationship between HPC and software engineering continues to evolve at a great pace.
It is important to understand cost and quality trade-offs as new technologies become available as well as the unique aspects of HPC software development.
